\ssr{ВВЕДЕНИЕ}

В настоящее время автоматическая проверка структуры текстовых документов является актуальной задачей,
особенно в образовательной сфере, где к оформлению отчётов, курсовых и выпускных квалификационных работ
предъявляются строгие требования. Одним из распространённых форматов хранения таких документов является PDF,
который затрудняет автоматический анализ структуры.

Целью данной лабораторной работы является разработка программного обеспечения на языке Python
для извлечения и анализа текстовых данных из PDF-файлов с использованием библиотеки PyPDF2 и регулярных выражений.
Вариант задания заключается в проверке корректности наличия структурных элементов
«ВВЕДЕНИЕ», «ЗАКЛЮЧЕНИЕ», «СПИСОК ИСПОЛЬЗОВАННЫХ ИСТОЧНИКОВ» в тексте документа и их присутствия в содержании,
а также в выявлении ошибочных наименований указанных элементов.

Для достижения поставленной цели в работе необходимо решить следующие задачи:
\begin{itemize}
    \item разработать регулярные выражения для проверки структурных элементов документа согласно варианту;
    \item реализовать функцию поиска подстрок в PDF-файле с использованием регулярных выражений;
    \item разработать программное обеспечение, принимающее на вход путь к PDF-файлу или каталогу с файлами;
    \item выполнить тестирование программы на наборе PDF-файлов, приложенных к лабораторной работе;
    \item провести сравнение решений, полученных с использованием различных больших языковых моделей.
\end{itemize}


